\section{Perspectives}
\label{sec:conclusion_perspectives}

The limitations discussed above open several directions for future research.
This section presents three of them, from the shortest to the longest term.

A first direction is to automate the semantic interpretation of the clusters.
One approach is to rely on vision-language models or \glspl{llm}.
Given observations from the same cluster, such models could generate a natural language description of what distinguishes it from the others.
This approach would provide a textual description for each cluster, helping identify its specific semantics and fully automate the \gls{car} method.

A second direction is to relate the identified clusters to the strategy followed by the agent.
The first step was to verify whether the latent space could be used to identify meaningful concepts.
The results obtained with \gls{car} show that this is a promising direction.
The next step is to use these concepts to explain the agent's strategy.
One possible approach would be to group similar states from the agent's perspective, thereby providing a simplified representation of the environment and capturing its relevant dynamics.

A third direction concerns the evaluation of explanations.
Structural metrics are necessary, but they are not sufficient.
The value of an explanation also depends on its effect on the observer.
Measuring this effect requires experimental protocols from disciplines such as cognitive science and psychology.
This would help determine whether the ultimate goal of explainability has been achieved.
